% Gradient boosting's first-order signal and XGBoost's second-order tree score.
% Schematic only; no empirical quantities are encoded.
% Canvas: 169 mm x 64 mm. Dependencies: project palette and TikZ.
\begin{tikzpicture}[
  foundation picture,x=1mm,y=1mm,
  font=\figurefont\footnotesize,
  state/.style={
    foundation node,
    draw=FigureBlue,
    fill=FigureBlueFill,
    minimum width=26mm,
    minimum height=10mm,
    text width=22mm,
    align=center,
    inner sep=1mm
  },
  model/.style={
    foundation node,
    draw=FigureRed,
    fill=FigureRedFill,
    minimum width=28mm,
    minimum height=10mm,
    text width=24mm,
    align=center,
    inner sep=1mm
  },
  signal/.style={
    foundation node,
    draw=FigureRed,
    fill=white,
    minimum width=30mm,
    minimum height=11mm,
    text width=26mm,
    align=center,
    inner sep=1mm
  },
  result/.style={
    foundation node,
    draw=FigureYellow,
    fill=FigureYellowFill,
    minimum width=29mm,
    minimum height=11mm,
    text width=25mm,
    align=center,
    inner sep=1mm
  },
  flow/.style={
    foundation flow
  }
]
  \path[use as bounding box] (0,0) rectangle (169,64);

  \node[anchor=west,foundation heading] at (1,60)
    {(a) 梯度提升：把损失方向变成新成员};
  \node[state] (current) at (18,42)
    {当前预测\\$F_{t-1}(\bm{x}_i)$};
  \node[signal] (gradient) at (52,42)
    {计算负梯度\\$r_i^{(t)}=-\partial_z\ell$};
  \node[model] (fit) at (52,19)
    {用$h_t$拟合\\伪残差};
  \node[result] (updated) at (18,19)
    {加法更新\\$F_t=F_{t-1}+\eta\alpha_t h_t$};
  \draw[flow] (current)--(gradient);
  \draw[flow] (gradient)--(fit);
  \draw[flow] (fit)--(updated);
  \draw[flow] (updated)--(current);

  \draw[FigureGrid,line width=0.6pt] (80,3)--(80,61);

  \node[anchor=west,foundation heading] at (85,60)
    {(b) XGBoost：二阶统计进入树评分};
  \node[signal] (derivatives) at (97,44)
    {样本统计\\$(g_i^{(t)},s_i^{(t)})$};
  \node[model] (partition) at (147,44)
    {候选分裂\\$I\to(I_{\mathrm L},I_{\mathrm R})$};
  \node[signal] (sums) at (99,22)
    {按候选叶聚合\\$(G_j,S_j)$};
  \node[result,minimum width=22mm,text width=18mm] (leaf) at (145,31)
    {叶权重\\$v_j^\star$};
  \node[result,minimum width=22mm,text width=18mm] (gain) at (145,13)
    {分裂增益\\$\operatorname{Gain}$};

  \draw[flow] (derivatives.east)--(partition.west);
  \draw[flow] (partition.south)--(147,35)--(99,35)--(sums.north);
  \draw[FigureInk,line width=.9pt] (sums.east)--(121,22);
  \fill[FigureInk] (121,22) circle[radius=1.1pt];
  \draw[flow] (121,22)|-(leaf.west);
  \draw[flow] (121,22)|-(gain.west);
\end{tikzpicture}
